\chapter{Arsitektur Sistem}
\label{ch:02-arsitektur}

\section{Gambaran Umum}

Sistem dibangun sebagai dua proses yang longgar terikat. Proses pertama adalah
\textbf{node perangkat}: satu sketch Arduino di \BoardName{} yang hanya mengenal
WiFi, broker MQTT, sensor ultrasonik, servo, modul RTC, dan LCD. Proses kedua adalah
\textbf{platform}: satu proses Python yang di dalamnya broker MQTT, worker ingesti,
sweeper, basis data SQLite, dan dashboard Streamlit hidup bersama. Keduanya tidak
pernah memanggil satu sama lain; satu-satunya jalur komunikasi adalah pesan MQTT
yang melewati broker. Konsekuensinya, perangkat tetap menjalankan jadwal harian dan
tetap menakar pakan ketika host mati, dan platform tetap menampilkan riwayat ketika
perangkat sedang tidak terpasang.

\figref{fig:e2e-topology} menempatkan kedua proses beserta berkas dan topik yang
menghubungkannya. Broker berada di dalam kotak proses Python: tidak ada layanan
MQTT yang perlu dipasang lebih dahulu.

\figureasset[\linewidth]{e2e-topology}{Topologi ujung-ke-ujung: node perangkat dan satu proses
Python yang memuat broker MQTT, worker, sweeper, basis data, dan dashboard.}{fig:e2e-topology}

\section{Blok Perangkat Lunak Platform}

Di dalam proses Python, tanggung jawab dipecah menjadi blok yang saling mengenal
hanya melalui basis data dan broker. \figref{fig:sw-block-diagram} menunjukkan
susunannya: satu entry point menyiapkan basis data, memilih mode broker, menyalakan
worker, lalu menyerahkan proses kepada Streamlit.

\figuretikz[\linewidth]{sw-block-diagram}{Blok perangkat lunak platform: entry point, pemilih
mode broker, worker ingesti, handler, sweeper, lapisan basis data, dan lapisan
presentasi.}{fig:sw-block-diagram}

Urutan penyalaan berlangsung satu arah. Pertama \code{db.init_db()} membuat skema
dan data awal bila tabel \code{devices} belum ada, lalu menjalankan migrasi kecil
yang idempoten. Kedua, platform memeriksa konfigurasi broker: bila
\code{EMBED_BROKER} bernilai salah, platform memakai broker pada \code{MQTT_URL};
bila benar, platform menguji apakah port sudah terpakai. Ketiga, worker paho
disambungkan secara asinkron dan diberi waktu lima detik untuk terhubung. Keempat,
sweeper dijalankan sebagai thread daemon. Kelima, proses diserahkan kepada Streamlit
pada port 8501. Urutan ini penting karena perangkat dapat menyambung kapan saja
tanpa menunggu ada orang membuka dashboard. (\file{dashboard/run.py},
\file{dashboard/runtime.py})

\section{Komponen Runtime}

\tabref{tab:02-komponen} memetakan berkas pada folder \file{dashboard} ke perannya.
Tabel ini menjadi rujukan pada bab berikutnya ketika sebuah perilaku perlu
ditelusuri ke implementasinya.

\begin{table}[H]
\centering
\caption{Komponen runtime platform dan perannya.}
\label{tab:02-komponen}
\small
\begin{tabular}{@{}p{2.4cm}p{12.2cm}@{}}
\toprule
\textbf{Berkas} & \textbf{Peran} \\
\midrule
\file{run.py} & Entry point: menyalakan runtime lebih dahulu, lalu menyerahkan
proses ke Streamlit sehingga broker dan worker sudah hidup sebelum ada sesi. \\
\file{runtime.py} & Singleton proses: menyiapkan basis data, memilih mode broker,
menyalakan worker, menjalankan sweeper, menyimpan label mode untuk sidebar. \\
\file{broker.py} & Broker MQTT tertanam (amqtt) dengan autentikasi anonim, plus
pemeriksa apakah port sudah dipakai proses lain. \\
\file{mqtt.py} & Worker paho: berlangganan empat jenis topik, memecah topik,
meneruskan payload ke handler, menyimpan tail pesan untuk halaman Live. \\
\file{handlers.py} & Empat handler pesan masuk, konversi jarak ke persen, dan
evaluasi alert. \\
\file{commands.py} & Jalur tulis keluar: menulis \code{commands} dan
\code{feeding_logs} sebelum mempublikasikan perintah. \\
\file{db.py} & Lapisan akses SQLite: koneksi thread-local, WAL,
\code{BEGIN IMMEDIATE}, seluruh kueri, backup dan restore. \\
\file{config.py} & Konfigurasi dari variabel lingkungan dan Streamlit Secrets. \\
\file{data.py} & Pembungkus baca untuk UI dengan cache berdurasi pendek. \\
\file{actions.py} & Aksi dari UI: mengirim perintah dan memberi umpan balik. \\
\file{app.py} & Kerangka dashboard: navigasi, pemilih perangkat, status sistem,
pemilih bahasa, panel backup dan restore. \\
\file{app_pages} & Isi tiap halaman: \file{status}, \file{live}, \file{alerts},
\file{history}, \file{schedules}, \file{control}. \\
\file{cli.py} & Pengiriman perintah tanpa dashboard lewat jalur tulis yang sama. \\
\bottomrule
\end{tabular}
\end{table}

\section{Pembagian Thread}

Satu proses Python memuat empat jalur eksekusi, seperti pada
\tabref{tab:02-thread}. Penulis basis data bukan hanya UI, sehingga setiap koneksi
bersifat thread-local dan setiap operasi multi-pernyataan dibungkus
\code{BEGIN IMMEDIATE}. (\file{dashboard/db.py})

\begin{table}[H]
\centering
\caption{Pembagian thread di dalam proses platform.}
\label{tab:02-thread}
\small
\begin{tabular}{@{}p{2.5cm}p{2.6cm}p{5.4cm}p{4.1cm}@{}}
\toprule
\textbf{Thread} & \textbf{Asal} & \textbf{Tugas} & \textbf{Sinkronisasi} \\
\midrule
main & Proses utama & Menjalankan skrip Streamlit per sesi, membaca basis data,
mengirim perintah. & Koneksi thread-local; cache baca pendek. \\
amqtt-broker & Daemon broker tertanam & Meneruskan pesan, menyimpan pesan retained,
menerbitkan LWT. & Event loop asyncio; tidak menyentuh basis data. \\
paho-network & Internal paho (\code{loop_start}) & Menjaga koneksi worker, memanggil
callback, menulis lewat handler. & Koneksi thread-local; transaksi untuk tulisan
berpasangan. \\
timeout-sweeper & Daemon yang dibuat runtime & Tiap \SweepS{} detik menandai perintah
tanpa ack dan perangkat tanpa data. & Transaksi singkat per sweep. \\
\bottomrule
\end{tabular}
\end{table}

Selain keempat thread itu, ada dua struktur memori bersama: tail pesan untuk
halaman Live berupa \code{deque} berukuran \LiveTail{} entri yang dilindungi
\code{threading.Lock}, dan singleton runtime yang dipakai seluruh sesi. State sesi
Streamlit tidak dipakai bersama antar thread karena hanya hidup di dalam sesi
peramban.

\section{Aliran Data Dua Arah}

Arah naik bersifat \emph{fire and forget}: telemetri dipublikasikan tiap
\TelemetryMs{} ms dengan QoS~0, status diterbitkan retained saat connect, dan event
jadwal diterbitkan setelah perangkat mengeksekusi jadwalnya sendiri. Arah turun
bersifat \emph{request-response}: platform menyimpan baris perintah lebih dahulu,
mempublikasikan dengan QoS~1, lalu menunggu ack yang membawa \code{command_id} yang
sama; bila ack tidak datang dalam \CmdTimeoutS{} detik, sweeper menyelesaikannya.

Pemisahan itu tercermin pada pekerjaan masing-masing pihak. Perangkat bertanggung
jawab atas aktuasi dan penjadwalan; platform atas pencatatan, agregasi, notifikasi,
dan penyajian. Persentase pakan dihitung kedua pihak dengan rumus yang sama, tetapi
platform selalu menghitung ulang dari \code{distance_cm} memakai kalibrasi pada
baris \code{devices}, sehingga kalibrasi dapat dikoreksi tanpa mem-flash ulang
firmware.

\section{Mengapa Broker Tertanam}

Broker tertanam dipilih karena tiga alasan praktis. Pertama, tidak ada
ketergantungan pada layanan eksternal: instalasi broker sistem operasi pernah
terblokir pada mesin pengembangan, dan aplikasi yang bergantung padanya tidak dapat
dijalankan sama sekali. Kedua, satu proses berarti satu perintah untuk menyalakan
seluruh platform. Ketiga, umur pakai lebih panjang karena broker hanya memerlukan
implementasi protokol yang benar, dan klien mana pun dapat dipakai untuk memeriksa
lalu lintas pesan.

Konsekuensinya, platform harus menangani tiga kemungkinan keadaan port seperti pada
\tabref{tab:02-mode-broker}. Mode terpilih ditampilkan di sidebar agar perilaku
sistem tidak menjadi misteri bagi pengguna. (\file{dashboard/broker.py})

\begin{table}[H]
\centering
\caption{Tiga mode broker yang mungkin terpilih saat platform dinyalakan.}
\label{tab:02-mode-broker}
\small
\begin{tabular}{@{}p{1.9cm}p{3.6cm}p{4.4cm}p{4.7cm}@{}}
\toprule
\textbf{Mode} & \textbf{Kapan terpilih} & \textbf{Perilaku} & \textbf{Catatan} \\
\midrule
embedded & \code{EMBED_BROKER} benar dan port \code{\BrokerBind} bebas & Broker amqtt
dinyalakan di dalam proses & Mode bawaan dan mode pengujian lokal. \\
existing & \code{EMBED_BROKER} benar tetapi port sudah dipakai & Platform tidak
menyalakan broker; worker menyambung ke \code{MQTT_URL} & Terjadi bila broker
sistem operasi berjalan sebagai layanan. \\
external & \code{EMBED_BROKER} salah & Platform menyambung ke broker pada
\code{MQTT_URL} & Dipakai untuk deployment yang tidak membuka port masuk. \\
\bottomrule
\end{tabular}
\end{table}

\begin{warnbox}
Mode \code{existing} dan \code{embedded} tidak dapat dibedakan dari sisi perangkat.
Karena itu perilaku broker tertanam yang menyimpang dari standar, yaitu ikut
mengirim pesan retained dari topik lain ke subscriber, harus ditangani di sisi
perangkat: firmware mengabaikan setiap pesan yang topiknya bukan \topic{command}
(\chapref{ch:06-protokol-mqtt}). Batas kepercayaan sistem ini adalah jaringan, bukan
aplikasi; akibatnya dicatat pada \chapref{ch:09-catatan-asumsi}.
\end{warnbox}
